%% Evaluation results for answers, deployment, verifier ordering and operations.
\subsection{E6: Answer and trace checks, abstention and cost}
\label{sec:e5}
\label{sec:agentic}\label{sec:abstention}\label{sec:composition}\label{sec:eval-agent}
Citation-marker counts differed across the three access paths in the original
grid and were equal in the repeat using the same five libraries. On ten research-paper
questions, the platform agent (P) produced nine answers with a
document-style marker, direct file tools (F) produced seven, and the
platform vector search over 800-character chunk copies (C) produced two.
In the repeat on the common libraries, each arm produced nine marker-bearing
answers (Figure~\ref{fig:answer}). Across cells matched by question identifier,
nine answers gained a marker and none lost one. BQ06 was worded differently
between the two grids, however, and the access substrates also changed; the
counts therefore do not identify a causal effect of substrate or question
wording. The chunk arm's BQ04 answer, for
instance, named source paths without using a marker recognised by the
check. Tables~S7--S8 of the supplementary inventory audit source presence,
while Tables~S14--S15 report the individual and aggregate grid checks.

The gold identifier appeared in the answer or trace in all 60 cells of two
30-cell grids (ten questions, three arms per grid).
None of the 60 saved question texts contains its corresponding gold ID, so
the retained question fields do not support input-ID copying as an explanation
for this result. The ID check searches answer text and recorded tool-call
inputs and agent text; it does not establish source resolution or semantic
support. The original answers were rescored using revised BQ06 annotations;
under the original annotation the original P arm scored 10/10 strict, not
9/10. Under the revised 60\% keyword threshold, original-grid strict
coverage was 9/10 in each arm, and alias-aware coverage was 10/10 in each.
On the common substrate, strict coverage was 7/10 for P, 7/10 for F and
8/10 for C, and alias-aware coverage
was 9/10, 9/10 and 10/10. These string checks do not establish semantic
correctness. In the original grid, the platform agent took 52.6~s and cost
\$0.711 per question. The corresponding means were 22.2~s and \$0.487 for
F, and 24.5~s and \$0.527 for C. A separate later run on the repaired
platform had different keyword and more permissive marker rules and is not
comparable to the historical grids. Its mean wall-clock time was 131~s per
question and its total cost was \$3.10 for ten questions.

\textbf{Answer-level case study.} After ingesting DeepSeek-R1, two agent
harnesses answered four questions with the library, paper and source identifier
provided and the retrieval, judging and reading sequence prescribed.
The case review checked answer content, source provenance, quotations
and fabricated facts against the paper and passed all eight responses
(Table~S16). On the electricity-cost question, both responses identified
the reported \$294K GPU-rental training estimate and reported that the
paper supplied no electricity figure. Three further industrial-corpus
probes returned not-found. Related absent probes still yielded retrieval
candidates above the threshold (Fig.~S3).

\begin{figure*}[!tb]
\centering
\includegraphics[width=\textwidth]{fig_opt_answer}
\caption{Document-style citation markers on two ten-question grids.
Each filled square represents one marker-bearing answer. (a) Original
grid with different access substrates. (b) Repeat on the same five
libraries. The gold identifier appeared in the answer or trace for
all 60 cells. The term-coverage counts are reported in Table~S15.
Each arm is a single execution.}
\label{fig:answer}
\end{figure*}
\begin{figure*}[!tb]
\centering
\includegraphics[width=\textwidth]{fig_opt_verifier}
\caption{Verifier scores and same-pool ordering from two separate
archives. (a) Laya scores for 191 recorded candidates across sixteen
hybrid runs. The 17 gold and 174 other scores have medians of 0.88 and
0.86, respectively. All recorded scores exceed 0.5, and one additional
Q14 candidate has no recorded score. Horizontal jitter separates the
marks. (b) Sixteen logged pools replayed in fusion and verifier order.
Filled marks indicate a gold document among the first five. Hit@5 is
15/16 in fusion order and 7/16 in verifier order.}
\label{fig:verifier}
\end{figure*}

\subsection{E7: Deployment-layer system arms and honesty contracts}
\label{sec:sysarms}
Three questions crossed pinned, unpinned and pinned hybrid-directed API
arms. The pinned arms returned both in-corpus gold identifiers, and all
three arms returned not-found for the out-of-corpus question. An unpinned
request exposed a transport limit. Its 131-document catalogue occupied
74.9~KB, exceeded the tool-output budget and yielded a truncated
preview, contributing to false refusals or stalled execution. Pagination and
grouped views let selected reruns complete, including the previously
failing query. A separate two-lane isolation test removed vector and
two-stage tools from one lane; its four initial traces recorded no vector-tool
calls, but none produced a final answer in that round. Catalogue truncation
was therefore not the only mechanism restricting access in these tests. In an
eighteen-arm battery, no final response contained a fabricated citation.
Interrupted arms were retried to produce one final response per condition.

\subsection{E8: When the verifier fails}
\label{sec:e6}
\label{sec:inflation}\label{sec:eval-verify}
On sixteen logged pools from the historical hybrid wrapper, ordering
candidates by verifier score reduced Hit@5 from 15/16 in fusion order to
7/16 (Figure~\ref{fig:verifier}b). Holding the recorded pools fixed isolates
the effect of sorting these candidates. The Q2 record illustrates the
demotion (Fig.~S5). The gold root-cause report scored 0.846, below five
telemetry reports scoring 0.887--0.926, and fell from first in fusion order
to eighth.

In the separate score archive, all 191 recorded values exceeded the
0.5 admission threshold. One additional Q14 candidate had no recorded
score. The medians were 0.88 for 17 gold candidates and 0.86 for 174
others (Figure~\ref{fig:verifier}a). The candidate-level Mann--Whitney
test ($p=0.26$) ignores clustering by query. Replays using an NLI
checkpoint and span scoring yielded Hit@5 of 0.3125 and 0.1875. A
relative-margin replay preserved 16 of 17 golds at $m=0.15$ and all 17
at $m=0.20$, retaining 10.4 and 11.1 candidates on average, respectively.
Table~S6 of the supplementary inventory and Fig.~S4 give the threshold
and margin sweeps. The paper corpus did not exhibit the same top-five
loss, with domain, language and corpus size changing together. The
replay identifies a risk of using admission scores as a ranking
criterion in this pool. Because all logged scores exceeded the configured
threshold, the archive provides no observed effect of rejecting
candidates at that threshold.

\subsection{E9: Operational boundaries: stability, concurrency, gate and
scope ablations}
\label{sec:e7}
The 27-cell stability run recorded success across three tasks, three lanes
and three rounds; these execution statuses do not establish that all 27
answers were semantically correct. In a separate
single-instance read-ladder concurrency sweep, the loads were 90, 180 and
270 requests at 5, 15 and 30 clients. The read-ladder recorded zero errors,
with median latency rising from 5.2 to 1577~ms. A separate vector-only
probe at 30 clients and 9 requests/s completed 159 of 270 requests within
30~s; the other 111 were not completed in that window, and are not
classified here as explicit timeout responses. A cache repair lowered the
30-client median to 65.8~ms on selected read-ladder reruns. These are two
pressure-test protocols, not one combined error/timeout result; the full
sweep after repair, concurrent writes and crash recovery were outside that
retest.

An inspection of cross-library vector search identified eight collections
with unreadable index segments. The global response skipped those
collections and reported success without exposing incomplete coverage.
Scoped search found a document in a newly ingested affected library,
and forced reindexing restored its global hit. A scoped probe on a
different affected library also returned no results. The implementation
logs per-collection errors and continues with the other collections, so
a completed response can still represent partial coverage.

A reconstructed fixed-candidate gate replay used ten questions, up to twelve
candidates and an evidence cap of 4{,}000 characters. All three arms
packed the gold document first on the eight answerable questions.
Filtered and unfiltered document sets matched on nine questions, while
LLM grades differed on some of these questions (including $0\to5$ and
$3\to0$). Matching document sets does not establish that the evidence-pack
text was identical, so these grade changes cannot be attributed solely to
generation or grading variability. Both questions without a gold document
were refused in every
arm. In the separate routing probes, A0 yielded two empty survivor
lists in four cases, both without a top candidate. A0h yielded none.
These outcomes distinguish absent candidates from rejection at the gate.

In a separate scope check, pinning the gold library yielded Hit@5 of
8/8 on the paper-corpus questions with in-corpus gold documents. Pinning a
non-gold library yielded 0/8. Catalogue routing selected the correct
library first for six of the eight. Filenames bore domain prefixes,
and the unscoped path searched 20 live libraries, including 15 beyond
the evaluation corpus.
